\documentclass[11pt, a4paper]{report}
%----------------------------------------------------------------------------------------
%	PACKAGES (宏包管理)
%----------------------------------------------------------------------------------------
% 1. 编码与语言
\usepackage[english]{babel}
\usepackage[T1]{fontenc}
% 2. 页面布局与排版
\usepackage{geometry}
\geometry{
    textwidth=6in, % 稍微加宽以便放入更好的box
    top=1in,
    bottom=1in,
    headheight=12pt,
    headsep=25pt,
    footskip=30pt
}
\usepackage{setspace}
\setstretch{1.2}
\usepackage{float}
\usepackage{enumitem} % 更好的列表控制
\usepackage{tocloft}  % 加载目录控制宏包
\renewcommand{\cftsecleader}{\cftdotfill{\cftdotsep}} % 强制让 section 后面显示点号
% 3. 数学支持
\usepackage{amsmath, amsthm, amssymb, amsfonts}
\usepackage{mathtools}
\usepackage{thmtools} % 定理环境增强
\usepackage{extarrows}
% 4. 图形与颜色
\usepackage{graphicx}
\usepackage[dvipsnames, svgnames]{xcolor} % 加载更多颜色名
\usepackage{tikz}
\usepackage{tikz-cd}
\usepackage{ytableau} % 杨表
\usetikzlibrary{decorations.pathreplacing, calc, shadows}

% Define reusable graph styles
\tikzset{
    vertex/.style={circle, fill=black, inner sep=1.5pt, outer sep=0pt},
    graphlabel/.style={align=center, font=\small}
}

% 5. 漂亮的文本框
\usepackage{tcolorbox}
\tcbuselibrary{breakable, skins, theorems}
% 6. 参考文献
%\usepackage[backend=biber,style=gb7714-2015,sorting=gb7714-2015]{biblatex}
%\addbibresource{references.bib}
% 7. 交互 (最后加载)
\usepackage[hidelinks]{hyperref}
\usepackage{array}
%----------------------------------------------------------------------------------------
%	CUSTOM COLORS (自定义配色方案)
%----------------------------------------------------------------------------------------
% 定理色系 - 蓝色
\definecolor{ThmColor}{RGB}{0, 71, 171}        % 深钴蓝
\definecolor{ThmBg}{RGB}{240, 244, 255}        % 极浅蓝
% 定义色系 - 绿色
\definecolor{DefColor}{RGB}{0, 100, 0}         % 深绿
\definecolor{DefBg}{RGB}{240, 255, 240}        % 蜜瓜绿
% 引理/命题色系 - 橙色
\definecolor{LemColor}{RGB}{204, 85, 0}        % 焦橙色
\definecolor{LemBg}{RGB}{255, 248, 240}        % 极浅橙
% 示例/注记色系 - 灰色/紫色
\definecolor{ExColor}{RGB}{105, 105, 105}      % 暗灰
\definecolor{ExBg}{RGB}{250, 250, 250}         % 极浅灰
% 专题/话题色系 - 靛紫色
\definecolor{TopicColor}{RGB}{75, 0, 130}      % 靛蓝色 (Indigo)
\definecolor{TopicBg}{RGB}{248, 245, 255}      % 极浅紫
%----------------------------------------------------------------------------------------
%	THEOREM ENVIRONMENTS (环境美化核心)
%----------------------------------------------------------------------------------------
% 1. Theorem (蓝色风格)
\declaretheorem[name=Theorem, numberwithin=section]{theorem}
\tcolorboxenvironment{theorem}{
    enhanced jigsaw,
    colframe=ThmColor,       % 边框颜色
    colback=ThmBg,           % 背景颜色
    coltitle=white,          
    boxrule=0.5pt,           % 边框粗细
    leftrule=2pt,            % 左边框加粗
    sharp corners=downhill,  % 样式
    breakable,               % 允许跨页
    drop shadow=black!5!white % 轻微阴影
}
% Corollary 使用与 Theorem 相同的风格，但通常紧随其后
\declaretheorem[name=Corollary, numberlike=theorem]{corollary}
\tcolorboxenvironment{corollary}{
    enhanced jigsaw,
    colframe=ThmColor!80!white, 
    colback=ThmBg,
    boxrule=0pt, leftrule=2pt, arc=0pt, outer arc=0pt,
    breakable
}
% 2. Definition (绿色风格)
\declaretheorem[name=Definition, numberlike=theorem]{definition}
\tcolorboxenvironment{definition}{
    enhanced jigsaw,
    colframe=DefColor,
    colback=DefBg,
    boxrule=0.5pt, leftrule=2pt,
    sharp corners=downhill,
    breakable,
    drop shadow=black!5!white
}
% 3. Lemma & Proposition (橙色风格)
\declaretheorem[name=Lemma, numberlike=theorem]{lemma}
\declaretheorem[name=Proposition, numberlike=theorem]{proposition}
\tcolorboxenvironment{lemma}{
    enhanced jigsaw,
    colframe=LemColor,
    colback=LemBg,
    boxrule=0.5pt, leftrule=2pt,
    breakable
}
\tcolorboxenvironment{proposition}{
    enhanced jigsaw,
    colframe=LemColor,
    colback=LemBg,
    boxrule=0.5pt, leftrule=2pt,
    breakable
}
% 4. Example & Remark (简约风格)
\declaretheorem[name=Example, numberlike=theorem]{example}
\declaretheorem[name=Remark, numberlike=theorem]{remark}
\declaretheorem[name=Exercise, numberlike=theorem]{exercise}
\tcolorboxenvironment{example}{
    enhanced jigsaw,
    colframe=ExColor,
    colback=ExBg,
    leftrule=2pt, boxrule=0pt,
    frame hidden, % 隐藏其他三边
    breakable
}
% 5. Topic (靛紫色风格，带标题支持)
\declaretheorem[name=Topic, numberlike=theorem]{topic}
\tcolorboxenvironment{topic}{
    enhanced jigsaw,
    breakable,
    colframe=TopicColor,        % 边框为靛紫色
    colback=TopicBg,           % 背景为极浅紫
    boxrule=0.5pt,             % 细边框
    leftrule=2pt,              % 左侧加粗线
    sharp corners=downhill,    % 与你其他环境一致的切角样式
    drop shadow=black!5!white, % 轻微阴影
    before skip=15pt,
    after skip=15pt
}
%----------------------------------------------------------------------------------------
%	MATH MACROS (保留你的常用定义)
%----------------------------------------------------------------------------------------
\newcommand{\g}{\mathfrak{g}}  
\newcommand{\der}{\mathrm{d}} 
\newcommand{\h}{\mathfrak{h}} 
\newcommand{\n}{\mathfrak{n}}  
\newcommand{\U}{\mathcal{U}}  
\newcommand{\C}{\mathbb{C}} 
\newcommand{\R}{\mathbb{R}} 
\newcommand{\borel}{\mathfrak{b}}
\newcommand{\Ind}{\mathrm{Ind}} 
\newcommand{\Hom}{\mathrm{Hom}}  
\newcommand{\Ext}{\mathrm{Ext}}
\newcommand{\HRule}[1]{\rule{\linewidth}{#1}}
%----------------------------------------------------------------------------------------
%	TITLE PAGE
%----------------------------------------------------------------------------------------
\begin{document}
\begin{titlepage}
    \centering
    \vspace*{2cm}
    
    \textsc{\Large Course Notes}\\[0.5cm]

    \HRule{1.5pt} \\[0.5cm]
    { \huge \bfseries Applied Abstract Algebra \\[0.3cm] Quiver Theory }\\[0.4cm] 
    \HRule{1.5pt} \\[1.5cm]
    
    \vspace{3cm}

    % 人员信息区域：垂直排列并区分字号
    \begin{spacing}{1.8}
        {\Large \textbf{Instructor:}} \\
        {\Large Prof. Iryna \textsc{Kashuba}} \\[0.8cm]
        
        {\normalsize \textit{TA:}} \\
        {\normalsize Leo}
    \end{spacing}

    \vfill
    {\large \today}
\end{titlepage}
\newpage
\tableofcontents
\newpage


%----------------------------------------------------------------------------------------
%    CONTENT STRUCTURE TEMPLATE
%----------------------------------------------------------------------------------------

\chapter{Lecture 1}

\begin{tcolorbox}[
    enhanced,
    title=\textbf{Course Information},
    colframe=TopicColor,
    colback=TopicBg,
    coltitle=white,
    boxrule=0.5pt,
    leftrule=2pt,
    sharp corners=downhill,
    drop shadow=black!5!white
]
\textbf{\textcolor{TopicColor}{Main topic:}} Representation theory of quivers.

\textbf{\textcolor{TopicColor}{Office hours:}} During Mon-Wed: 9:00-16:00, we can discuss for 2 hours of office hours.

\textbf{\textcolor{TopicColor}{Main book:}} \textit{Ralf Schiffler, Quiver Representations}, Cambridge University Press.

\textbf{\textcolor{TopicColor}{References:}}
\begin{itemize}[leftmargin=*, nosep, label=\textcolor{TopicColor}{$\bullet$}]
    \item \textit{Assem, Simons, Skowronski, Elements of the Representation Theory of Associative Algebras}, Cambridge University Press.
    \item \textit{Coelho, modules, algebras and quivers}.
    \item \textit{Michel Brion, Representations of quivers}.
    \item \textit{Henning Krause, Representations of quivers via reflection functor}.
    \item \textit{W.Crawley-Boevey, Lectures on representations of quivers}.
    \item \textit{W.Crawley-Boevey, More on representations of quivers}.
    \item Other Reference: \url{www.math.uni-bielefeld.de}
\end{itemize}
\end{tcolorbox}

\section{Introduction}

\begin{definition} 
A quiver $Q$ is a quadruple $(Q_0, Q_1, s, t)$, where $Q_0$ is a set of vertices, $Q_1$ is a set of arrows, and $s, t: Q_1 \to Q_0$ are two maps. For any arrow $\alpha \in Q_1$, $s(\alpha)$ is called the source of $\alpha$ and $t(\alpha)$ is called the target of $\alpha$. For example, if $ \alpha: 1 \to 2 $, then $ s(\alpha)=1 $ and $ t(\alpha)=2 $.
\end{definition}

\begin{example} 
Let $ Q_0=\{ 1,2,3 \} $, $ Q_1=\{ \alpha,\beta, \gamma, \delta \} $.
We have $ s(\alpha)=2, t(\alpha)=2 $; $ s(\beta)=3, t(\beta)=1 $; $ s(\gamma)=s(\delta)=1, t(\gamma)=t(\delta)=2 $.
\begin{center}
\begin{tikzcd}
\underset{\bullet}{1} \arrow[r, "\gamma", bend left] \arrow[r, "\delta"', bend right] & \underset{\bullet}{2} \arrow["\alpha"', loop, distance=2em, in=305, out=235] & \underset{\bullet}{3} \arrow[ll, "\beta"', bend right=49]
\end{tikzcd}
\end{center}
Let path $ \xi=\alpha \gamma \beta $, then we have $ s(\xi)=s(\beta)=3 $ and $ t(\xi)=t(\alpha)=2 $.  
\end{example}

\begin{example} 
Jordan quiver.
The Jordan quiver consists of a single vertex and a single arrow starting and ending at this vertex. Formally, $Q_0 = \{1\}$ and $Q_1 = \{\alpha\}$ with $s(\alpha) = t(\alpha) = 1$.
It is called the Jordan quiver because its representations correspond to vector spaces equipped with a linear endomorphism. The classification of such representations up to isomorphism is equivalent to the classification of square matrices up to similarity, which is given by the Jordan normal form theorem (over an algebraically closed field).
\begin{center}
\begin{tikzcd}
\bullet \arrow[loop, distance=2em, in=125, out=55]
\end{tikzcd}
\end{center}
\end{example}

\begin{example} 
$ r $-Kronecker quiver.
The $r$-Kronecker quiver has two vertices, say 1 and 2, and $r$ arrows $\alpha_1, \dots, \alpha_r$ from 1 to 2. Thus $Q_0 = \{1, 2\}$, $Q_1 = \{\alpha_1, \dots, \alpha_r\}$, with $s(\alpha_i) = 1$ and $t(\alpha_i) = 2$ for all $i=1, \dots, r$.
It is named after Leopold Kronecker. In the case $r=2$, the representations correspond to pairs of linear maps between two vector spaces, which can be represented by matrix pencils. Kronecker solved the classification problem for these pencils.
\begin{center}
    \begin{tikzcd}
\bullet \arrow[r, "1", bend left=49] \arrow[r, "\ldots", no head, dashed] \arrow[r, "r"', bend right=49] & \bullet
\end{tikzcd}
\end{center}
\end{example}

\begin{example} 
$ n $-subspace quiver. 
The $n$-subspace quiver consists of $n+1$ vertices $\{0, 1, \dots, n\}$ and $n$ arrows $\alpha_1, \dots, \alpha_n$ such that $s(\alpha_i) = i$ and $t(\alpha_i) = 0$ for all $i=1, \dots, n$.
It is called the $n$-subspace quiver because its representations are related to configurations of subspaces. If the linear maps associated with the arrows are injective, a representation corresponds to a collection of $n$ subspaces of the vector space associated with the central vertex.
\begin{center}
\begin{tikzcd}
                                  & \underset{\bullet}{0}            &        &                                    \\
                                  &                                  &        &                                    \\
\underset{\bullet}{1} \arrow[ruu] & \underset{\bullet}{2} \arrow[uu] & \ldots & \underset{\bullet}{n} \arrow[lluu]
\end{tikzcd}
\end{center}
\end{example}

\begin{example} 
Finite quivers.
A quiver $Q=(Q_0, Q_1, s, t)$ is called finite if both the set of vertices $Q_0$ and the set of arrows $Q_1$ are finite sets.
\end{example}

For the base field we assume $ k=\bar{k} $. Usually we take $ k=\C $. 

\begin{definition} 
A representation $M$ of a quiver $Q$ consists of a vector space $M_i$ for each vertex $i \in Q_0$ and a linear map $\phi_{\alpha}: M_{s(\alpha)} \to M_{t(\alpha)}$ for each arrow $\alpha \in Q_1$. We denote this representation by $M=(M_i, \phi_{\alpha})_{i \in Q_0, \alpha \in Q_1}$.
\end{definition}

\begin{example} 
For example 1: 
\begin{center}
\begin{tikzcd}
\underset{\bullet}{M_1} \arrow[r, "\phi_\gamma", bend left] \arrow[r, "\phi_\delta"', bend right] & \underset{\bullet}{M_2} \arrow["\phi_\alpha"', loop, distance=2em, in=305, out=235] & \underset{\bullet}{M_3} \arrow[ll, "\phi_\beta"', bend right=49]
\end{tikzcd}
\end{center}
\end{example}

\begin{definition} 
A representation $M$ is called \textbf{finite dimensional} if $\dim M_i < \infty$ for all $i \in Q_0$. The tuple $\dim M = \{\dim M_i\}_{i \in Q_0}$ is called the \textbf{dimension vector} of $M$. An element $m \in M$ is a tuple $\{m_i\}_{i \in Q_0}$ where $m_i \in M_i$ for each $i \in Q_0$.
\end{definition}

\begin{example} 
$ 1 $-Kronecker quiver:
\begin{center}
    \begin{tikzcd}
1 \arrow[r] & 2
\end{tikzcd}
\end{center}

Take $ k $ to be a field. Some representations of this quiver are:
\[
M:\;
\begin{tikzcd}[ampersand replacement=\&, column sep=large]
k \arrow[r, "1"] \& k
\end{tikzcd}
\qquad
M':\;
\begin{tikzcd}[ampersand replacement=\&, column sep=large]
k \arrow[r, "0"] \& k
\end{tikzcd}
\qquad
M'':\;
\begin{tikzcd}[ampersand replacement=\&, column sep=large]
k \arrow[r, "0"] \& 0
\end{tikzcd}
\]

Another example is
\[
\widetilde{M}:\;
\begin{tikzcd}[ampersand replacement=\&, column sep=huge]
k^2 \arrow[r, "{\begin{bmatrix}1 & 0 \\ 1 & 0 \\ 0 & 0\end{bmatrix}}"] \& k^3
\end{tikzcd}.
\]
For $ v=(1,0) \in k^2 $, we have
\[
\begin{bmatrix}
1 & 0 \\
1 & 0 \\
0 & 0
\end{bmatrix}
\begin{bmatrix}
1 \\
0
\end{bmatrix}
=
\begin{bmatrix}
1 \\
1 \\
0
\end{bmatrix}.
\]
\end{example}

\begin{definition}[Morphism between two representations of $ Q $] 
Let $ M=(M_i, \phi_{\alpha})_{i \in Q_0, \alpha \in Q_1} $ and $ N=(N_i, \psi_{\alpha})_{i \in Q_0, \alpha \in Q_1} $ be two representations of $ Q $. A morphism $ f: M \to N $ is a collection of linear maps $ f_i: M_i \to N_i $ such that for every $ \phi_{\alpha}: M_i \to M_j $, $ \alpha \in Q_1 $ the following diagram commutes:
    \begin{center}
    \begin{tikzcd}
    M_i \arrow[r, "\phi_{\alpha}"] \arrow[d, "f_i"'] & M_j \arrow[d, "f_j"] \\
    N_i \arrow[r, "\psi_{\alpha}"']                  & N_j                 
    \end{tikzcd}
    \end{center}
\end{definition}

\begin{example} 
In example of 1-Kronecker quiver, let $ f: M \to M'' $, $ f_1: k \to k $ such that $ f_1(1)=a $, and $ f_2: k \to 0 $ is the zero map. We can check that $ f $ is a morphism. Check the following diagram:
\begin{center}
    \begin{tikzcd}
k \arrow[r, "1"] \arrow[d, "f_1 \ \ \ a"'] & k \arrow[d, "0 \ \ \ f_2"] \\
k \arrow[r, "0"']                          & 0                         
\end{tikzcd}
\end{center}

Now we establish $ g: M'' \to M $. If $ g $ is a morphism, then we have $ g_1: k \to k $ such that $ g_1(1)=b $, and $ g_2: 0 \to k $ is the zero map. To make the target diagram commute it forces $ g_1=0 $. Hence there are only $ 0 $ morphism form $ M'' $ to $ M $.
\begin{center}
\begin{tikzcd}
k \arrow[r, "0"] \arrow[d, "g_1 \ \ \ b"'] & 0 \arrow[d, "0 \ \ \ g_2"] \\
k \arrow[r, "1"']                          & 0                         
\end{tikzcd}
\end{center}

\[ 1 \circ g_1 = 0 \Rightarrow g_1=0 \]
\end{example}

Recall the definition of category:

\begin{definition} 
A category $\mathcal{C}$ consists of the following data:
\begin{itemize}
    \item A class of objects, denoted by $\mathrm{Ob}(\mathcal{C})$.
    \item For any pair of objects $X, Y \in \mathrm{Ob}(\mathcal{C})$, a set of morphisms from $X$ to $Y$, denoted by $\mathrm{Hom}_{\mathcal{C}}(X, Y)$. If $f \in \mathrm{Hom}_{\mathcal{C}}(X, Y)$, we write $f: X \to Y$.
    \item For any three objects $X, Y, Z \in \mathrm{Ob}(\mathcal{C})$, a composition map:
    \[
    \circ: \mathrm{Hom}_{\mathcal{C}}(Y, Z) \times \mathrm{Hom}_{\mathcal{C}}(X, Y) \to \mathrm{Hom}_{\mathcal{C}}(X, Z), \quad (g, f) \mapsto g \circ f.
    \]
\end{itemize}
These data must satisfy the following axioms:
\begin{enumerate}
    \item \textbf{Associativity}: For any morphisms $f: X \to Y$, $g: Y \to Z$, and $h: Z \to W$, we have
    \[
    h \circ (g \circ f) = (h \circ g) \circ f.
    \]
    \item \textbf{Identity}: For every object $X \in \mathrm{Ob}(\mathcal{C})$, there exists an identity morphism $\mathrm{id}_X \in \mathrm{Hom}_{\mathcal{C}}(X, X)$ such that for any $f: X \to Y$ and $g: Z \to X$,
    \[
    f \circ \mathrm{id}_X = f \quad \text{and} \quad \mathrm{id}_X \circ g = g.
    \]
\end{enumerate}
\end{definition}

Denote $ \mathrm{Rep}Q $ the category of representations of $ Q $. Then $ \mathrm{Ob}(\mathrm{Rep}Q) $ is the class of representations of $ Q $: $M=(M_i, \phi_{\alpha})$, and $ \mathrm{Hom}_{\mathrm{Rep}Q} $ is the class of morphisms $  \mathrm{Hom}_{\mathrm{Rep}Q}(M,N) $ from $ M $ to $ N $.

Let $ M,N \in \mathrm{Ob}(\mathrm{Rep}Q) $. Then $ \mathrm{Hom}(M,N):= \mathrm{Hom}_{\mathrm{Rep}Q}(M,N) $ is a vector space over $k$ with operations defined component-wise: for $ f,g \in \mathrm{Hom}(M,N) $ and $ \lambda \in k $, let $(f+g)_i = f_i + g_i$ and $(\lambda f)_i = \lambda f_i$ for all $ i \in Q_0 $.

\textbf{Verification:} Let $ M=(M_i, \phi_\alpha) $ and $ N=(N_i, \psi_\alpha) $. Since $ f,g $ are morphisms, for any arrow $ \alpha: i \to j $, we have $ \psi_\alpha f_i = f_j \phi_\alpha $ and $ \psi_\alpha g_i = g_j \phi_\alpha $.
\begin{enumerate}
    \item \textbf{Addition:}
    \[
    \psi_\alpha (f+g)_i = \psi_\alpha (f_i + g_i) = \psi_\alpha f_i + \psi_\alpha g_i = f_j \phi_\alpha + g_j \phi_\alpha = (f_j + g_j) \phi_\alpha = (f+g)_j \phi_\alpha.
    \]
    Thus $ f+g \in \mathrm{Hom}(M,N) $.
    \item \textbf{Scalar Multiplication:}
    \[
    \psi_\alpha (\lambda f)_i = \psi_\alpha (\lambda f_i) = \lambda (\psi_\alpha f_i) = \lambda (f_j \phi_\alpha) = (\lambda f_j) \phi_\alpha = (\lambda f)_j \phi_\alpha.
    \]
    Thus $ \lambda f \in \mathrm{Hom}(M,N) $.
\end{enumerate}
The zero morphism $ 0=(0)_{i \in Q_0} $ clearly satisfies the commutativity condition. The vector space axioms follow from the structure of linear maps between vector spaces.

\begin{example} 
Consider 
 \begin{center}
    \begin{tikzcd}
1 \arrow[r] & 2
\end{tikzcd}
\end{center}
$ \mathrm{Hom}(M,M'') =\{ (a,0) | a \in k \} \simeq k $. $\mathrm{Hom}(M'',M) =\{ 0\} $. 
\end{example}

\begin{example} 
Let $ Q $ be a quiver:
\begin{center}
\begin{tikzcd}
1 \arrow[rd, "\alpha"] &   &                        \\
                       & 3 & 4 \arrow[l, "\gamma"'] \\
2 \arrow[ru, "\beta"]  &   &                       
\end{tikzcd}
\end{center}
Let $ M $:
\begin{center}
    \begin{tikzcd}
k \arrow[rd, "\begin{bmatrix} 1 \\ 0 \end{bmatrix}"] &     &                                                      \\
                                                     & k^2 & k \arrow[l, "\begin{bmatrix} 1 \\ 1 \end{bmatrix}"'] \\
k \arrow[ru, "\begin{bmatrix} 0 \\ 1 \end{bmatrix}"] &     &                                                     
\end{tikzcd}
\end{center}

$ M' $:
\begin{center}
\begin{tikzcd}
0 \arrow[rd, "0"] &   &                   \\
                  & k & 0 \arrow[l, "0"'] \\
0 \arrow[ru, "0"] &   &                  
\end{tikzcd}
\end{center}

$ M'' $:
\begin{center}
\begin{tikzcd}
k \arrow[rd, "1"] &   &                   \\
                  & k & k \arrow[l, "1"'] \\
k \arrow[ru, "1"] &   &                  
\end{tikzcd}
\end{center}

Let us check $ \mathrm{Hom}(M,M') $:
\begin{center}
\begin{tikzcd}
k \arrow[rd, "\begin{bmatrix} 1 \\ 0 \end{bmatrix}"] \arrow[ddd, "0"', bend right] &                                &                                                                                  \\
                                                                                   & k^2 \arrow[ddd, "{[a \ \ b]}"] & k \arrow[l, "\begin{bmatrix} 1 \\ 1 \end{bmatrix}"'] \arrow[ddd, "0", bend left] \\
k \arrow[ru, "\begin{bmatrix} 0 \\ 1 \end{bmatrix}"] \arrow[ddd, "0"', bend right] &                                &                                                                                  \\
0 \arrow[rd, "0"]                                                                  &                                &                                                                                  \\
                                                                                   & k                              & 0 \arrow[l, "0"]                                                                 \\
0 \arrow[ru, "0"]                                                                  &                                &                                                                                 
\end{tikzcd}
\end{center}

Since $M'_1=M'_2=M'_4=0$, the maps $f_1, f_2, f_4$ are zero. Let $f_3 = \begin{bmatrix} a & b \end{bmatrix}$. The commutativity conditions impose:
\[
\begin{cases}
f_3 \phi_\alpha = \psi_\alpha f_1 \implies \begin{bmatrix} a & b \end{bmatrix} \begin{bmatrix} 1 \\ 0 \end{bmatrix} = 0 \implies a = 0, \\
f_3 \phi_\beta = \psi_\beta f_2 \implies \begin{bmatrix} a & b \end{bmatrix} \begin{bmatrix} 0 \\ 1 \end{bmatrix} = 0 \implies b = 0.
\end{cases}
\]
Thus $f_3 = 0$, and $\mathrm{Hom}(M, M') = 0$.

Let us check $ \mathrm{Hom}(M,M'') $:
\begin{center}
\begin{tikzcd}
k \arrow[rd, "\begin{bmatrix} 1 \\ 0 \end{bmatrix}"] \arrow[ddd, "c"', bend right] &                                &                                                                                  \\
                                                                                   & k^2 \arrow[ddd, "{[a \ \ b]}"] & k \arrow[l, "\begin{bmatrix} 1 \\ 1 \end{bmatrix}"'] \arrow[ddd, "e", bend left] \\
k \arrow[ru, "\begin{bmatrix} 0 \\ 1 \end{bmatrix}"] \arrow[ddd, "d"', bend right] &                                &                                                                                  \\
0 \arrow[rd, "0"]                                                                  &                                &                                                                                  \\
                                                                                   & k                              & 0 \arrow[l, "0"]                                                                 \\
0 \arrow[ru, "0"]                                                                  &                                &                                                                                 
\end{tikzcd}
\end{center}
Let the morphism be given by $g_1 = [c]$, $g_2 = [d]$, $g_4 = [e]$, and $g_3 = \begin{bmatrix} a & b \end{bmatrix}$. The commutativity conditions are:
\[
\begin{cases}
g_3 \phi_\alpha = \psi_\alpha g_1 \implies \begin{bmatrix} a & b \end{bmatrix} \begin{bmatrix} 1 \\ 0 \end{bmatrix} = c \implies a = c, \\
g_3 \phi_\beta = \psi_\beta g_2 \implies \begin{bmatrix} a & b \end{bmatrix} \begin{bmatrix} 0 \\ 1 \end{bmatrix} = d \implies b = d, \\
g_3 \phi_\gamma = \psi_\gamma g_4 \implies \begin{bmatrix} a & b \end{bmatrix} \begin{bmatrix} 1 \\ 1 \end{bmatrix} = e \implies a + b = e.
\end{cases}
\]
The morphism is uniquely determined by the parameters $a, b \in k$. Hence $\mathrm{Hom}(M, M'') \cong k^2$.
\end{example}

\begin{example} 
Dynkin Diagrams:

\begin{center}
$A_l \ (l \ge 1)$:\quad
\begin{tikzpicture}[baseline=-0.5ex]
\fill (0,0) circle (2pt);
\fill (1.5,0) circle (2pt);
\fill (3.0,0) circle (2pt);
\fill (6.0,0) circle (2pt);
\fill (7.5,0) circle (2pt);
\draw (0,0)--(1.5,0)--(3.0,0);
\draw (6.0,0)--(7.5,0);
\node at (4.5,0.35) {$\cdots$};
\node[below] at (0,0) {$1$};
\node[below] at (1.5,0) {$2$};
\node[below] at (3.0,0) {$3$};
\node[below] at (6.0,0) {$l-1$};
\node[below] at (7.5,0) {$l$};
\end{tikzpicture}
\end{center}

\begin{center}
$D_l \ (l \ge 4)$:\quad
\begin{tikzpicture}[baseline=-0.5ex]
\fill (0,0) circle (2pt);
\fill (1.5,0) circle (2pt);
\fill (4.5,0) circle (2pt);
\fill (6.0,0) circle (2pt);
\fill (8.2,1.5) circle (2pt);
\fill (8.2,-1.5) circle (2pt);
\draw (0,0)--(1.5,0);
\node at (3.0,0.35) {$\cdots$};
\draw (4.5,0)--(6.0,0);
\draw (6.0,0)--(8.2,1.5);
\draw (6.0,0)--(8.2,-1.5);
\node[below] at (0,0) {$1$};
\node[below] at (1.5,0) {$2$};
\node[below] at (4.5,0) {$l-3$};
\node[below] at (6.0,0) {$l-2$};
\node[right] at (8.2,1.5) {$l-1$};
\node[right] at (8.2,-1.5) {$l$};
\end{tikzpicture}
\end{center}

\begin{center}
$E_6$:\quad
\begin{tikzpicture}[baseline=-0.5ex]
\foreach \x in {0,1.5,3.0,4.5,6.0}{
  \fill (\x,0) circle (2pt);
}
\fill (3.0,1.5) circle (2pt);
\draw (0,0)--(1.5,0)--(3.0,0)--(4.5,0)--(6.0,0);
\draw (3.0,0)--(3.0,1.5);
\node[below] at (0,0) {$1$};
\node[below] at (1.5,0) {$3$};
\node[below] at (3.0,0) {$4$};
\node[below] at (4.5,0) {$5$};
\node[below] at (6.0,0) {$6$};
\node[above] at (3.0,1.5) {$2$};
\end{tikzpicture}
\end{center}

\begin{center}
$E_7$:\quad
\begin{tikzpicture}[baseline=-0.5ex]
\foreach \x in {0,1.5,3.0,4.5,6.0,7.5}{
  \fill (\x,0) circle (2pt);
}
\fill (3.0,1.5) circle (2pt);
\draw (0,0)--(1.5,0)--(3.0,0)--(4.5,0)--(6.0,0)--(7.5,0);
\draw (3.0,0)--(3.0,1.5);
\node[below] at (0,0) {$1$};
\node[below] at (1.5,0) {$3$};
\node[below] at (3.0,0) {$4$};
\node[below] at (4.5,0) {$5$};
\node[below] at (6.0,0) {$6$};
\node[below] at (7.5,0) {$7$};
\node[above] at (3.0,1.5) {$2$};
\end{tikzpicture}
\end{center}

\begin{center}
$E_8$:\quad
\begin{tikzpicture}[baseline=-0.5ex]
\foreach \x in {0,1.5,3.0,4.5,6.0,7.5,9.0}{
  \fill (\x,0) circle (2pt);
}
\fill (3.0,1.5) circle (2pt);
\draw (0,0)--(1.5,0)--(3.0,0)--(4.5,0)--(6.0,0)--(7.5,0)--(9.0,0);
\draw (3.0,0)--(3.0,1.5);
\node[below] at (0,0) {$1$};
\node[below] at (1.5,0) {$3$};
\node[below] at (3.0,0) {$4$};
\node[below] at (4.5,0) {$5$};
\node[below] at (6.0,0) {$6$};
\node[below] at (7.5,0) {$7$};
\node[below] at (9.0,0) {$8$};
\node[above] at (3.0,1.5) {$2$};
\end{tikzpicture}
\end{center}
\end{example}

\begin{example} 
$ 2 $-Kronecker quiver:
\begin{center}
    \begin{tikzcd}
1 \arrow[r, bend left] \arrow[r, bend right] & 2
\end{tikzcd}
\end{center}
Representations of this quiver:
\begin{center}
    \begin{tikzcd}
V \arrow[r, bend left, "A"] \arrow[r, bend right, "B"'] & W
\end{tikzcd}
\end{center}
\end{example}

\begin{example} 
Jordan quiver:
\begin{center}
\begin{tikzcd}
1 \arrow[loop, distance=2em, in=125, out=55]
\end{tikzcd}
\end{center}
Representations of this quiver:
\begin{center}
\begin{tikzcd}
V \arrow["A"', loop, distance=2em, in=125, out=55]
\end{tikzcd}
\end{center}
$ A:V\to V $, the matrix representation of $ A $ is similar to a Jordan normal form (when $ k=\bar{k} $).  
\end{example}

\chapter{Lecture 2}

\begin{definition}
Let $Q=(Q_0,Q_1)$ be a quiver. A representation of $Q$ is a collection
\[
M=(V_i,\varphi_\alpha)_{i \in Q_0,\alpha \in Q_1},
\]
where each $V_i$ is a vector space and each arrow $\alpha:i \to j$ is assigned a linear map
\[
\varphi_\alpha \in \operatorname{Hom}(V_i,V_j).
\]
\end{definition}

\begin{example}
Consider the $2$-Kronecker quiver
\[
\begin{tikzcd}
1 \arrow[r, bend left] \arrow[r, bend right] & 2
\end{tikzcd}
\]
and the representations
\[
M:\quad
k \xrightarrow{\left(\begin{bmatrix}1\\0\end{bmatrix},\,\begin{bmatrix}0\\1\end{bmatrix}\right)} k^2
\]
and
\[
N:\quad
k^2 \xrightarrow{\left(\begin{bmatrix}1&1\\0&1\end{bmatrix},\,\begin{bmatrix}1&1\\1&1\end{bmatrix}\right)} k^2.
\]
Let $f=(f_1,f_2):M \to N$ be a morphism. Write
\[
f_1=\begin{bmatrix}a\\b\end{bmatrix},
\qquad
f_2=\begin{bmatrix}c&d\\e&f\end{bmatrix}.
\]
Commutativity along the two arrows gives
\[
f_2\begin{bmatrix}1\\0\end{bmatrix}
=
\begin{bmatrix}c\\e\end{bmatrix}
=
\begin{bmatrix}1&1\\0&1\end{bmatrix}
\begin{bmatrix}a\\b\end{bmatrix}
=
\begin{bmatrix}a+b\\b\end{bmatrix},
\]
and
\[
f_2\begin{bmatrix}0\\1\end{bmatrix}
=
\begin{bmatrix}d\\f\end{bmatrix}
=
\begin{bmatrix}1&1\\1&1\end{bmatrix}
\begin{bmatrix}a\\b\end{bmatrix}
=
\begin{bmatrix}a+b\\a+b\end{bmatrix}.
\]
Hence
\[
c=d=f=a+b,
\qquad
e=b,
\]
so
\[
f=
\left(
\begin{bmatrix}a\\b\end{bmatrix},
\begin{bmatrix}a+b&a+b\\b&a+b\end{bmatrix}
\right),
\qquad
\dim \operatorname{Hom}(M,N)=2.
\]
This illustrates how morphisms in $\operatorname{Rep}(Q)$ are determined by commutative diagrams.
\end{example}

If $M=(V_i,\varphi_\alpha)$ and $N=(U_i,\psi_\alpha)$ are representations of $Q$, then their direct sum is
\[
M \oplus N=
\left(
V_i \oplus U_i,
\begin{bmatrix}
\varphi_\alpha & 0 \\
0 & \psi_\alpha
\end{bmatrix}
\right).
\]
Inductively, one may form finite direct sums $M_1 \oplus \cdots \oplus M_r$.

\begin{definition}
A representation $M$ is called indecomposable if it cannot be written as a direct sum of two non-zero representations of $Q$.
\end{definition}

\begin{remark}
For quiver representations, indecomposable and irreducible are different notions; they do not have to coincide.
\end{remark}

\begin{example}
For the Jordan quiver, a representation is a pair $(V,A)$ with $A:V \to V$. Two such representations are isomorphic if and only if the corresponding matrices are similar. Therefore the isomorphism classes are controlled by Jordan normal form, and the indecomposable representations correspond exactly to single Jordan blocks.
\end{example}

\begin{theorem}[Krull--Schmidt]
Every finite-dimensional representation in $\operatorname{Rep}(Q)$ can be decomposed as a finite direct sum of indecomposable representations. Moreover, this decomposition is unique up to reordering and isomorphism of the summands.
\end{theorem}



\chapter{Lecture 3}

\begin{example} 
\begin{tikzcd}
1 \arrow[r] & 2
\end{tikzcd} We have trivial representation 
\[
\begin{tikzcd}
0 \arrow[r, "0"] & 0
\end{tikzcd}
\]
The following 3 representations are non-trivial
\[
\begin{tikzcd}
k \arrow[r, "0"] & 0
\end{tikzcd}
\begin{tikzcd}
0 \arrow[r, "0"] & k
\end{tikzcd}
\begin{tikzcd}
k \arrow[r, "1"] & k
\end{tikzcd}
\]
Denoted by $ S_1, S_2, M $.

First two representations are irriducible and indecomposable, the last one is decomposable. Recall that for a morophism $ f: M \to N $, we have induced representation 
\[
\begin{tikzcd}
\operatorname{ker} f \arrow[r, "incl.", hook] & M
\end{tikzcd}
\]
And a subrepresentation of $ M $ is representation $ L $ such that $ f: L \to M $ is a monomorphism. 

Now consider
\[
\begin{tikzcd}
k \arrow[r, "0"] \arrow[d, "a"'] & 0 \arrow[d, "b"] \\
k \arrow[r, "1"]                 & k               
\end{tikzcd}
\]
Then a=b=0, it is impossible to construct inclusion.

But if we consider
\[
\begin{tikzcd}
0 \arrow[r, "0"] \arrow[d, "a"'] & k \arrow[d, "b"] \\
k \arrow[r, "1"]                 & k               
\end{tikzcd}
\]
Then $ a=b \cdot 0=0 $ and $ b $ is arbitrary (then we have a monomorphism). Hence we come to the conclusion that $ \operatorname{Hom}(S_1,M)=0 $ and $ \operatorname{Hom}(S_2,M)=k $. In other words $ S_2 $ is a subrepresentation of $ M $ and hence $ M $ is reducible.

But if $ S_2 $ is a direct summand of $ M $, then consider   



\end{example}

\begin{definition}[equivalent definition of $ \operatorname{ker}f$] 

Let $ f:M\to N $ be a morphism of representations. We call $ \operatorname{ker}f $ a representation $ g: L \to M $ such that $ fg=0 $ such that for any honomorphism $ h:X\to M $ with $ fh=0 $, there exists a unique morphism $ u:X\to L $ such that $ g u = h $. In other words, the following diagram commutes:

\begin{center}
\begin{tikzcd}
X \arrow[r, "u", dashed] \arrow[rd, "h"'] & L \arrow[d, "g"] \\
& M \arrow[r, "f"] & N
\end{tikzcd}
\end{center}
\end{definition}

Let us show that $ \operatorname{ker}f=(\operatorname{ker}f_i=L_i, \varphi_{\alpha}|_{L_i}) $ is equivalent to the diagram representation $ L=\operatorname{ker}f , g_i=ind.$ i.e. $ \begin{tikzcd}
\operatorname{ker}f_i \arrow[r, "incl.", hook] & M_i
\end{tikzcd} $ where $ f_ig_i=0 $.

Let $ h: X\to M $ be a morphism such that $ fh=0 $. Let $ u: X\to L $. Consider $ (M_i , \varphi_\alpha) , (L_i, \psi_\alpha)$ and $ (X_i, \nu_\alpha) $. $\forall  x_i \in X_i $, $ f_ih_i(x_i)=0 $ i.e. $ h_i(x_i) \in \operatorname{ker}f_i =L_i $. Define $ u_i(x_i)=h_i(x_i) $.

We have to check that $ u $ is a morphism and $ h=gu $. Consider the following diagram:
  \[
  \begin{tikzcd}
X_i \arrow[r, "\nu_\alpha"] \arrow[d, "u_i"'] & X_j \arrow[d, "u_j"] \\
L_i \arrow[r, "\psi_\alpha"]                  & L_j                 
\end{tikzcd}
  \]

Analogously we have equivalent definition of $ \operatorname{coker}f $.

\begin{definition} 
Let $ f : M\to N $ be a morphism of representations. $ f: f_i \in \operatorname{Hom}(M_i, N_i)_{i \in Q_0} $. $ \operatorname{coker}f =(N_i/ f_i(M_i), \chi_\alpha(n_i+f_i(M_i))=\psi_\alpha (n_i), f_j(M_j)) $. Equivalently, $ \operatorname{coker}f $ is a representation $ g: N\to H $ such that $ gf=0 $ and for any morphism $ h: N \to Y $ with $ hf=0 $, there exists a unique morphism $ u: H \to Y $ such that $ ug=h $. In other words, the following diagram commutes:  
\end{definition}

\begin{exercise}
Prove the quivalence of the two definitions of $ \operatorname{coker}f $.
\end{exercise}

\begin{theorem}[firsi isomorphism theorem] 
If $ f: M\to N $ is a morphism of representations, then $ \operatorname{im}f \cong M/\operatorname{ker}f $. 
\end{theorem}

\begin{proof} 
$ (\operatorname{im}f)_i = \operatorname{im}f_i \subseteq N_i $ and $ \operatorname{im}f =(\operatorname{im}f_i, \psi_\alpha |_{\operatorname{im}f_i}) $. Check that this is a representation. Let $ (M/\operatorname{ker}f)_i = M_i / \operatorname{ker}f_i $, $ M/\operatorname{ker}f =(M_i / \operatorname{ker}f_i, \overline{\varphi_{\alpha}}) $ where $\overline{\varphi_{\alpha}}  $ is defined as $ \overline{\varphi_{\alpha}}(m_i + \operatorname{ker})=\varphi_\alpha(m_i)+ \operatorname{ker}f_j $.

From the isomorphism theorem for vector spaces, $ \bar{f_i}(m_i + \operatorname{ker}f_i) := f_i(m_i) $, $ \bar{f_i}: M_i/\operatorname{ker}f_i \to \operatorname{im}f_i $. We need to check that $ \bar{f} : M / \operatorname{ker}f \cong \operatorname{im}f$.   

Consider the following diagram:
\[
\begin{tikzcd}
M_i / \operatorname{ker}f_i \arrow[r, "\overline{\varphi_{\alpha}}"] \arrow[d, "u_i"'] & M_j / \operatorname{ker}f_j \arrow[d, "u_j"] \\
\operatorname{im}f_i \arrow[r, "\psi_\alpha |_{\operatorname{im}f_i}"]                 & \operatorname{im}f_j                        
\end{tikzcd}
\]
Check the commutativity of the diagram:
\end{proof}

These all together shows that $ \operatorname{Rep}Q $ is abelian category. More precisely,

\begin{definition} 
We call a category $ \mathcal{C} $ abelian $ k $-category if the following conditions are satisfied:
\begin{enumerate}
    \item $\mathcal{C} $ is $ k $-categor, for any $ X,Y \in \mathrm{Ob}(\mathcal{C}) $, $ \operatorname{Hom}_{\mathcal{C}}(X,Y) $ is a vector space over $ k $.
    \item $ \mathcal{C} $ is a additive category, $ \mathcal{C} $ has direct sums and zero object.
    \item Each morphism $ f: X \to Y $ has kernel $ K $ and cokernel $ H $. $ f: M\to N $, $ i : K\to M $, $ p: N\to H $. and $ \operatorname{coker}i \cong \operatorname{ker}p $.    
\end{enumerate}
\end{definition}

\begin{definition} 
Let $ \begin{tikzcd}
M \arrow[r] & N \arrow[r] & P
\end{tikzcd} $ be a sequence of morphisms in $ \operatorname{Rep}Q $. We say that this sequence is exact if $ \operatorname{ker}p = \operatorname{im}f $.
\end{definition}

\begin{definition} 
We call a sequence of morphisms in $ \operatorname{Rep}Q $a short exact sequence if it is exact and of the form
\[
\begin{tikzcd}
0 \arrow[r] & L \arrow[r, "f"] & M \arrow[r, "g"] & N \arrow[r] & 0
\end{tikzcd}
\]
Where $ f $ is injective, $ g $ is surjective $ \operatorname{im}f = \operatorname{ker}g $  
\end{definition}

For $ f: M\to N $ we have exact sequence:
\[
\begin{tikzcd}
0 \arrow[r] & \operatorname{ker}f \arrow[r, "i"] & M \arrow[r, "f"] & N \arrow[r, "p"] & \operatorname{coker}f \arrow[r] & 0
\end{tikzcd}
\] 
Where $ i $ is an inclusion and $ p $ is a projection.

Or we have short exact sequence:
\[
\begin{tikzcd}
0 \arrow[r] & \operatorname{ker}f \arrow[r, "i"] & M \arrow[r, "p"] & M/\operatorname{ker}f \arrow[r] & 0
\end{tikzcd}
\]

\begin{example} 
picture
\end{example}

\begin{definition} 
A morphism $ f: L\to M $ a section if there exists a morphism $ h: M\to L $ such that $ hf=1_L $. A morphism $ g: M\to N $ is a retraction if there exists a morphism $ h: N\to M $ such that $ gh=1_N $. 
\end{definition}

\begin{definition} 
We say that a short exact sequence 
\[ \begin{tikzcd}
0 \arrow[r] & L \arrow[r, "f"] & M \arrow[r, "g"] & N \arrow[r] & 0
\end{tikzcd} 
\]
splits if $ f $ has a retraction or equivalently $ f $ is a section.   
\end{definition}

\begin{theorem} 
Let 
\[
\begin{tikzcd}
0 \arrow[r] & L \arrow[r, "f"] & M \arrow[r, "g"] & N \arrow[r] & 0
\end{tikzcd}
\]
be a short exact sequence in $ \operatorname{Rep}Q $. Then 
\begin{enumerate}
    \item $ f $ is a section if and only if $ g $ is a retraction.
    \item If $ f $ is a seciton then $ \operatorname{ker}g = \operatorname{im}f $ is a direct summand of $ M $.
\end{enumerate}
\end{theorem}

\begin{proof} 
$ (\Rightarrow): $ Suppose $ f $ is a section. $ \exists h : M\to L $ such that $ hf=1_L $. We will define $ h': N\to M $ and show that $ gh'=1_N $. $ \forall n \in N $, since $ g $ is surjective $ \exists m $ such that $ n =g(m) $. $ h'(n)= m - f(h(m)) $ where $ h': N\to M $. We have to show that $ h' $ is well-defined. 

Suppose $ m,m' $ such that $ n=g(m)=g(m') $. We have to show that $ m'-f(h(m'))=m-f(h(m)) $. $ m'-m  =f(h(m'))-f(h(m)) \in \operatorname{im}f$, $ \exists l \in L $ such that $ f(l)=m'-m= $. Now:
\[
m' -f(h(m')) -m + f(h(m))= m'-m-f(h(m'-m))=m'-m-f(h(f(l)))=m'-m-f(l)=0
\]                
\end{proof}






\end{document}
